\section{The topological balance and the moving wall}
\label{sec:boundary}

We orient $\C^2$ by $dt\wedge ds\wedge dx\wedge dy$.
All boundaries below are piecewise smooth, with the outward
orientation. Corners have zero three-dimensional measure.

\begin{lemma}[Local heat charge]
\label{lem:charge}
Suppose $H\not\equiv0$ is holomorphic near a point $p=(\tau_*,z_*)$,
satisfies $H_\tau=-H_{zz}$, and $H(\tau_*,\cdot)$ has a zero
of order $n\ge2$ at $z_*$. Then the joint zero of $(H,H_z)$ is
isolated, and its local complex degree is $\binom n2$.
\end{lemma}
\begin{proof}
Translate $p$ to zero and write the leading spatial coefficient as
$a\ne0$. The convergent heat expansion has weighted leading term
\[
 H(\tau,z)=aP_n(\tau,z)+\text{terms of weighted degree greater than }n,
 \qquad
 P_n=\sum_{j=0}^{\lfloor n/2\rfloor}
       \frac{n!(-\tau)^jz^{n-2j}}{j!(n-2j)!},
\]
where $\tau$ has weight two and $z$ weight one.
The polynomial $P_n(1,z)$ is a scaled Hermite polynomial and has
$n$ distinct real roots. One elementary proof uses its orthogonality
with a positive Gaussian weight: fewer than $n$ distinct sign changes
would produce a polynomial of degree below $n$ whose product with
$P_n$ has one sign, contradicting orthogonality.

Put $\tau=v^2$ and $z=vw$. The quotient $v^{-n}H(v^2,vw)$
extends holomorphically to $v=0$ and equals $aP_n(1,w)$ there.
The implicit function theorem gives all $n$ roots in the cluster as
$z_j(v^2)=v(r_j+O(v))$, with distinct $r_j$.
Consequently its monic Weierstrass discriminant is
\[
 \Delta(\tau)=c\tau^{n(n-1)/2}(1+o(1)),\qquad c\ne0.
\]
It is holomorphic in $\tau$, so its zero has the indicated integer
order and is isolated. This also isolates the joint zero.

For completeness, the local intersection number of $H=0$ and
$H_z=0$ equals the vanishing order of this discriminant. Indeed,
Weierstrass preparation writes $H=UP$, with $U$ a unit and $P$
monic of degree $n$. Modulo $P$, $H_z=UP_z$; multiplication by
the unit changes the resultant by a nonvanishing factor. The
resultant of $P$ and $P_z$ is its discriminant up to sign. Equivalently,
the determinant of multiplication by $P_z$ on the free
$\C\{\tau\}$-module $\C\{\tau\}[z]/(P)$ has this vanishing
order, equal to the dimension of the local intersection algebra.
A small regular perturbation counts that many positive holomorphic
intersections, hence gives the local degree.
\end{proof}

At real coefficients the same expansion shows that, just after a real
event, all $n$ local roots are simple and real. Just before it there
are $n\bmod2$ real roots. Thus the real birth weight and complex
degree stated in Section~\ref{sec:intro} follow. Positivity of the
degree is not a sign estimate on each individual boundary face.

\begin{lemma}[Complete boundary integral]
\label{lem:degree}
Let $W$ be a compact domain on whose boundary $G$ is nonzero.
Then
\begin{equation}
 \int_{\partial W}\CC
 =\deg(G/|G|\,;\partial W)
 =\sum_{p\in W:\,G(p)=0}\binom{m_p}{2}.
 \label{eq:degree}
\end{equation}
\end{lemma}
\begin{proof}
On the unit sphere in $\C^2$, the form
$\alpha=\im(\bar\zeta_1d\zeta_1+\bar\zeta_2d\zeta_2)$ satisfies
$\alpha\wedge d\alpha=2\,d\mathrm{Vol}_{S^3}$.
Its integral is $4\pi^2$. This follows by evaluating at $(1,0)$
on $(\partial_{y_1},\partial_{x_2},\partial_{y_2})$ and using
unitary invariance. The form $\CC$ is the pullback of this normalized
volume. Excise disjoint balls around the finitely many joint zeros
and apply degree additivity, then Lemma~\ref{lem:charge}.
\end{proof}

\subsection{The exact cap integral}

A regular projective root tube over a complex-time rectangle $D$
is a holomorphic inverse
\begin{equation}
 Z:D\times\{|q|\le\epsilon\}\longrightarrow\C,
 \qquad H_z(\tau,Z)\ne0,\qquad
 \frac{H(\tau,Z)}{H_z(\tau,Z)}=q,
 \label{eq:inverse}
\end{equation}
defined on a neighborhood of the displayed compact set and injective
in each $q$ fiber. The radius is constant on this whole tube.
Differentiating the inverse equation gives $Z_q\ne0$, so its
parameterization preserves the complex spatial orientation. Its
center $z(\tau)=Z(\tau,0)$ follows a simple zero.

\begin{proposition}[Phase reduction and zero side flux]
\label{prop:phase}
For $m$ disjoint tubes of the form \eqref{eq:inverse}, with a common
radius $\epsilon$, define
\begin{equation}
 h_j(\tau)=H_z(\tau,z_j(\tau)),\qquad
 \DD(\tau)=\prod_{j=1}^m h_j(\tau).
 \label{eq:D}
\end{equation}
For an oriented heat edge $e$, the integral over its root disks is
\begin{equation}
 I(e)=\frac{k_\epsilon}{2\pi}\im\int_e d\log\DD.
 \label{eq:phase}
\end{equation}
Every wall $|q|=\epsilon$ has identically zero pullback of $\CC$.
Furthermore,
\begin{equation}
 \DD=(-1)^{m(m-1)/2}\Delta\UU.
 \label{eq:factor}
\end{equation}
The factor $U=H/\prod_j(z-z_j)$ need only be a nonvanishing
holomorphic factor near the selected roots. It can have zeros outside
those neighborhoods; no infinite product is required.
\end{proposition}
\begin{proof}
In one tube write $h(\tau,q)=H_z(\tau,Z(\tau,q))$ and
$\phi=\im\log h$, using a continuous logarithm on the contractible
parameter domain. The physical jet is exactly $h(q,1)$, so
\begin{align}
 \Aconn&=d\phi+a(q),&
 a(q)&=\frac{\im(\bar q\,dq)}{1+|q|^2},\nonumber\\
 d\Aconn&=\omega(q),&
 \omega(q)&=\frac{2\,du\wedge dv}{(1+u^2+v^2)^2},\quad q=u+iv.
 \label{eq:chart}
\end{align}
Since $a\wedge\omega=0$,
$\CC=d\phi\wedge\omega/(4\pi^2)$. The derivative of $\phi$
in the edge direction is harmonic in $q$. Its weighted average over
the disk is its center value times
\[
 \int_{|q|\le\epsilon}\omega
 =\int_0^{2\pi}\int_0^\epsilon\frac{2r}{(1+r^2)^2}\dd r\dd\vartheta
 =2\pi k_\epsilon.
\]
This proves \eqref{eq:phase} after summing the tubes.
On the wall $q=\epsilon e^{i\vartheta}$, the two-form $\omega$
pulls back to zero, proving the side assertion.

At a selected simple root,
$H_z(\tau,z_j)=U(\tau,z_j)\prod_{i\ne j}(z_j-z_i)$.
Multiplying these identities pairs the factors belonging to every
unordered pair and gives \eqref{eq:factor}.
\end{proof}

The vanishing side flux is a geometric statement, not a bound on the
speed of a root. It remains true when the disks move rapidly. A
radius allowed to vary \emph{along} the tube would produce
$dr\wedge d\vartheta$ terms and would not have the same property.
We allow a different constant radius for each chosen terminal $T$.

\begin{theorem}[The complete transport formula]
\label{thm:transport}
Suppose $W$ is an event-containing compact domain with $G\ne0$ on
its boundary, whose terminal face is the union of $m$ disks from
Proposition~\ref{prop:phase}. Decompose its other boundary into the
initial face, the two imaginary-time faces, regular projective walls,
and all remaining spatial or joining faces. Their oriented integrals
are $-\II_0$, $\JJ_+$, $\JJ_-$, $0$, and $\RR$, respectively.
Then \eqref{eq:organizing} holds with \eqref{eq:terms}.

If $A(W)$ is the number of distinct real creation events in the
interior, and separately proved bounds give
\[
 \II_0\ge L_0,\quad \II_T\le U_T,\quad
 \JJ_++\JJ_-\le U_J,\quad \RR\le U_R,
\]
then
\begin{equation}
 A(W)\le\QQ\le U_T+U_J+U_R-L_0.
 \label{eq:eventbound}
\end{equation}
\end{theorem}
\begin{proof}
Lemma~\ref{lem:degree} gives
$\QQ=\II_T-\II_0+\JJ_++\JJ_-+\RR$.
Substitute Proposition~\ref{prop:phase}. Each real event contributes
at least one to the nonnegative sum defining $\QQ$, proving the bound.
\end{proof}

This is a useful interface for estimates, but it does not prove that
the four required one-sided bounds have favorable sizes. In
particular, a bound for a closure expressed in terms of $\QQ$ itself
characterizes the closure; it cannot be inserted into
\eqref{eq:eventbound} as independent information about $\QQ$.

\subsection{The other two faces of a regular tube}

Assume real type, meaning
$H(\bar\tau,\bar z)=\overline{H(\tau,z)}$, and let the root branches be real at real time on
$[a,b]$. Each $h_j(t)$ then has a constant real sign. Choose
\begin{equation}
 \Theta(t,Y)=\sum_j\im\log\frac{h_j(t+iY)}{h_j(t)},
 \label{eq:Theta}
\end{equation}
where each logarithm is continued vertically from zero imaginary
time. Conjugation and Proposition~\ref{prop:phase} give
\begin{equation}
 \boxed{\displaystyle
 \II_t=\frac{k_\epsilon}{\pi}\Theta(t,Y),\qquad
 \JJ_+=\JJ_-=
 \frac{k_\epsilon}{2\pi}\{\Theta(a,Y)-\Theta(b,Y)\}
 =\frac{\II_a-\II_b}{2}.}
 \label{eq:closures}
\end{equation}
The orientation on the upper imaginary face is
$-dt\wedge du\wedge dv$; on the lower it is
$dt\wedge du\wedge dv$. Thus the complete regular tube obeys
$\II_b-\II_a+\JJ_++\JJ_-=0$, as it must because it contains
no joint zero. Attaching tubes to an event region does not remove
the complement of their attachment disks.
